\documentclass[10pt,twocolumn]{article}

% ---------------------------------------------------------------
% Packages
% ---------------------------------------------------------------
\usepackage[margin=1in]{geometry}
\usepackage{amsmath,amssymb,amsthm}
\usepackage{graphicx}
\usepackage{hyperref}
\usepackage{url}
\usepackage{booktabs}
\usepackage{listings}
\usepackage{xcolor}
\usepackage{microtype}
\usepackage{natbib}
\usepackage{algorithm}
\usepackage{algpseudocode}
\usepackage{enumitem}
\usepackage{caption}
\usepackage{subcaption}

% ---------------------------------------------------------------
% Code listing style (PyTorch / Python)
% ---------------------------------------------------------------
\definecolor{codebg}{HTML}{F7F4EE}
\definecolor{codegreen}{rgb}{0,0.5,0}
\definecolor{codegray}{rgb}{0.45,0.45,0.45}
\definecolor{codeblue}{HTML}{2d5f8a}
\definecolor{codepurple}{HTML}{7a3e9d}

\lstset{
  backgroundcolor=\color{codebg},
  basicstyle=\ttfamily\scriptsize,
  breaklines=true,
  captionpos=b,
  commentstyle=\color{codegreen}\itshape,
  keywordstyle=\color{codeblue}\bfseries,
  stringstyle=\color{codepurple},
  language=Python,
  frame=single,
  framesep=4pt,
  showstringspaces=false,
  tabsize=4,
  xleftmargin=8pt,
  xrightmargin=4pt,
  numbers=none,
}

% ---------------------------------------------------------------
% Theorem environments
% ---------------------------------------------------------------
\newtheorem{theorem}{Theorem}
\newtheorem{proposition}[theorem]{Proposition}
\newtheorem{definition}{Definition}
\newtheorem{remark}{Remark}

% ---------------------------------------------------------------
% Hyperref setup
% ---------------------------------------------------------------
\hypersetup{
  colorlinks=true,
  linkcolor=codeblue,
  citecolor=codeblue,
  urlcolor=codeblue,
}

% ---------------------------------------------------------------
% Title
% ---------------------------------------------------------------
\title{Open Questions in AI}
\author{
  Mohan Chinnappan \\
  \texttt{mohan.chinnappan.n@gmail.com}
}

\date{June 2026}

% ---------------------------------------------------------------
\begin{document}
% ---------------------------------------------------------------

\maketitle

% ---------------------------------------------------------------
\begin{abstract}
% ---------------------------------------------------------------
Modern large language models (LLMs) achieve remarkable performance on a
broad range of tasks, yet we lack a principled theory explaining
\emph{why} they reason, generalise, or hallucinate the way they do.
This paper surveys five grand open questions at the frontier of AI
research: (1)~\textbf{mechanistic interpretability} — can we reverse-engineer
neural networks into human-understandable circuits?; (2)~\textbf{world models
and causality} — do LLMs learn implicit causal models, and can explicit
world models do better?; (3)~\textbf{reasoning and computational depth} — are
transformer reasoning abilities genuine, and what are their fundamental
limits?; (4)~\textbf{representation collapse and JEPA} — does predicting in
latent space rather than signal space produce more structured, interpretable
representations?; and (5)~\textbf{modular AI} — will future general systems
be monolithic or composed of specialised coordinated modules?  For each
question we provide precise mathematical formulations, discuss the current
state of knowledge, and illustrate key concepts with PyTorch code.  We
argue that answering these questions requires not only scaling, but
fundamentally new architectural inductive biases.
\end{abstract}

\textbf{Keywords:} interpretability, world models, JEPA, mechanistic
interpretability, causal reasoning, sparse autoencoders, self-supervised
learning, energy-based models.

% ---------------------------------------------------------------
\section{Introduction}
% ---------------------------------------------------------------

The past five years have witnessed extraordinary advances in AI: LLMs
such as GPT-4 \citep{openai2023gpt4}, Gemini \citep{gemini2023}, and
Claude \citep{anthropic2024claude} solve complex reasoning problems,
write code, and pass professional examinations.  Despite this empirical
success, we lack a complete theoretical understanding of \emph{why these
systems work} and \emph{where they fundamentally fail}.

This gap has produced entire research sub-fields: \emph{mechanistic
interpretability} attempts to reverse-engineer neural network circuits
\citep{olah2020zoom}; \emph{representation learning} studies the geometry of
learned feature spaces \citep{elhage2022superposition}; \emph{world model
research} asks whether models build causal models of their environments
\citep{lecun2022path}; and \emph{scaling law research} asks whether the
current paradigm can reach general intelligence \citep{hoffmann2022chinchilla}.

We organise the open questions as follows.  Section~\ref{sec:interpretability}
covers mechanistic interpretability.  Section~\ref{sec:worldmodels} covers
world models and causality.  Section~\ref{sec:reasoning} covers reasoning
and computational limits.  Section~\ref{sec:jepa} covers the JEPA
architecture as a concrete alternative inductive bias.
Section~\ref{sec:modular} discusses the modular AI hypothesis.
Section~\ref{sec:conclusion} concludes with an integrated research agenda.

% ---------------------------------------------------------------
\section{Background: What an LLM Computes}
\label{sec:background}
% ---------------------------------------------------------------

An autoregressive language model parameterised by $\theta$ computes

\begin{equation}
  P_\theta(y_1, \ldots, y_T) = \prod_{t=1}^{T} P_\theta(y_t \mid y_{<t})
  \label{eq:lm}
\end{equation}

\noindent where $y_{<t} = (y_1, \ldots, y_{t-1})$ is the context.  Training
minimises the negative log-likelihood:

\begin{equation}
  \mathcal{L}_{\text{NLL}}(\theta) = -\frac{1}{T}\sum_{t=1}^{T} \log P_\theta(y_t \mid y_{<t}).
  \label{eq:nll}
\end{equation}

A transformer \citep{vaswani2017attention} implements this through $L$
alternating self-attention and feed-forward layers.  Each self-attention
head computes

\begin{equation}
  \text{Attn}(Q,K,V) = \text{softmax}\!\left(\frac{QK^\top}{\sqrt{d_k}}\right)V,
  \label{eq:attention}
\end{equation}

\noindent where $Q = h W_Q$, $K = h W_K$, $V = h W_V$ are linear projections
of the hidden state $h \in \mathbb{R}^{d_{\text{model}}}$.  The full
residual-stream update is

\begin{align}
  h_\ell &= \text{LayerNorm}(h_{\ell-1} + \text{MHA}(h_{\ell-1})) \label{eq:res1}\\
  h_\ell'&= \text{LayerNorm}(h_\ell + \text{FFN}(h_\ell)).\label{eq:res2}
\end{align}

\noindent Critically, Equation~\ref{eq:nll} contains no term that rewards
physical understanding, causal reasoning, or goal-directed planning.
Any such capability must \emph{emerge} from next-token prediction alone —
and this is the root cause of all five open questions.

\subsection{PyTorch: Basic Transformer Block}

\begin{lstlisting}[caption={Single transformer layer in PyTorch.}]
import torch, torch.nn as nn, math

class TransformerBlock(nn.Module):
    def __init__(self, d_model, n_heads, d_ff):
        super().__init__()
        self.attn = nn.MultiheadAttention(
            d_model, n_heads, batch_first=True)
        self.ff = nn.Sequential(
            nn.Linear(d_model, d_ff),
            nn.GELU(),
            nn.Linear(d_ff, d_model))
        self.ln1 = nn.LayerNorm(d_model)
        self.ln2 = nn.LayerNorm(d_model)

    def forward(self, x):
        # x: (B, T, d_model)
        h = self.ln1(x)
        h, _ = self.attn(h, h, h)
        x = x + h                      # Eq. (3)
        x = x + self.ff(self.ln2(x))   # Eq. (4)
        return x
\end{lstlisting}

% ---------------------------------------------------------------
\section{Open Question 1: Mechanistic Interpretability}
\label{sec:interpretability}
% ---------------------------------------------------------------

\begin{definition}[Mechanistic Interpretability]
The program of \emph{mechanistic interpretability} aims to identify
human-understandable algorithms implemented by the weights of a neural
network — analogous to reverse-engineering a compiled binary into
readable source code \citep{olah2020zoom}.
\end{definition}

\subsection{The Superposition Hypothesis}

A key obstacle is \emph{polysemanticity}: single neurons respond to
multiple, semantically unrelated concepts.  \citet{elhage2022superposition}
propose the \textbf{superposition hypothesis}: a network of dimension
$d$ can represent up to $n \gg d$ features by encoding them as
near-orthogonal directions in $\mathbb{R}^d$, at the cost of
inter-feature interference.

\begin{proposition}[Superposition Capacity, informal]
If features are sparse — each is active with probability $p \ll 1$ —
then $n \approx d \cdot (-\log p)^{-1}$ features can be stored with
interference below a threshold $\epsilon$.
\end{proposition}

The activation of a polysemantic neuron $j$ for input $\mathbf{x}$ is

\begin{equation}
  h_j = \text{ReLU}\!\left(\sum_{i} W_{ji} x_i + b_j\right),
  \label{eq:neuron}
\end{equation}

\noindent where multiple features $i$ share weight $W_{ji}$, making
clean attribution impossible.

\subsection{Sparse Autoencoders}

Sparse autoencoders (SAEs) \citep{bricken2023monosemanticity} decompose
polysemantic activations into monosemantic features by solving

\begin{align}
  \mathbf{z} &= \text{ReLU}(W_e \mathbf{h} + \mathbf{b}_e)
                \label{eq:sae_enc}\\
  \hat{\mathbf{h}} &= W_d \mathbf{z} + \mathbf{b}_d
                \label{eq:sae_dec}\\
  \mathcal{L}_{\text{SAE}} &= \|\mathbf{h} - \hat{\mathbf{h}}\|_2^2
                + \lambda \|\mathbf{z}\|_1,
                \label{eq:sae_loss}
\end{align}

\noindent where $W_d \in \mathbb{R}^{d_{\text{model}} \times d_{\text{SAE}}}$
with $d_{\text{SAE}} \gg d_{\text{model}}$.  Each column of $W_d$ is a
candidate \emph{feature direction}.

\begin{lstlisting}[caption={Sparse autoencoder for mechanistic interpretability.}]
class SparseAutoencoder(nn.Module):
    def __init__(self, d_model, d_sae):
        super().__init__()
        self.W_enc = nn.Linear(d_model, d_sae)
        self.W_dec = nn.Linear(d_sae, d_model,
                               bias=False)
        self.b_enc = nn.Parameter(
            torch.zeros(d_sae))

    def forward(self, h):
        # Eq. (5)
        z = torch.relu(self.W_enc(h) + self.b_enc)
        h_hat = self.W_dec(z)           # Eq. (6)
        return h_hat, z

def sae_loss(h, h_hat, z, lam=1e-3):
    recon = (h - h_hat).pow(2).mean()  # Eq. (7)
    sparse = lam * z.abs().mean()
    return recon + sparse
\end{lstlisting}

\subsection{Linear Representation Hypothesis}

The \textbf{linear representation hypothesis} posits that concepts are
encoded as linear directions in activation space.  Formally, concept
$c$ is encoded at layer $\ell$ if there exists $\mathbf{v}_c \in \mathbb{R}^d$
such that

\begin{equation}
  \hat{c}(\mathbf{h}_\ell) = \text{sign}(\mathbf{v}_c^\top \mathbf{h}_\ell - \theta)
  \label{eq:probe}
\end{equation}

\noindent achieves high accuracy.  Causal evidence requires that patching

\begin{equation}
  \mathbf{h}_\ell \leftarrow \mathbf{h}_\ell + \alpha \mathbf{v}_c
  \label{eq:patch}
\end{equation}

\noindent changes the model's output in the direction predicted by concept $c$.

\subsection{The Scaling Barrier}

\begin{remark}
Despite impressive progress, mechanistic interpretability faces a
fundamental scaling barrier: manually annotating circuits for a model
with $\sim 10^{12}$ parameters is computationally infeasible.
Automated methods (SAEs, probing) cover only a small fraction of
model behaviour, and there is no ground-truth evaluation criterion
for whether an interpretation is ``correct.''
\end{remark}

% ---------------------------------------------------------------
\section{Open Question 2: World Models and Causality}
\label{sec:worldmodels}
% ---------------------------------------------------------------

\begin{definition}[World Model]
A world model is a parametric function $f_\theta : \mathcal{S} \times
\mathcal{A} \to \mathcal{S}$ that predicts future states from current
states and actions, trained to minimise

\begin{equation}
  \mathcal{L}_{\text{WM}}(\theta) = \mathbb{E}_{(s_t, a_t, s_{t+1})}
  \!\left[\|f_\theta(s_t, a_t) - s_{t+1}\|_2^2\right].
  \label{eq:wm}
\end{equation}
\end{definition}

\subsection{The Blurry Prediction Problem}

When the target $s_{t+1}$ is stochastic and the prediction is made
directly in signal space (pixels, tokens), the optimal predictor is
the conditional mean

\begin{equation}
  \hat{s}_{t+1}^* = \mathbb{E}[s_{t+1} \mid s_t, a_t],
  \label{eq:cond_mean}
\end{equation}

\noindent which is a blurry average of all possible futures.  If a ball
can move left or right with equal probability, the optimal
$\hat{s}^*$ is a ghost ball at the centre — physically impossible.

\subsection{Causal Reasoning}

Pearl's \citeyear{pearl2009causality} \textbf{do-calculus} distinguishes
observational from interventional distributions:

\begin{align}
  \text{Observational:} &\quad P(Y \mid X = x) \label{eq:obs}\\
  \text{Interventional:} &\quad P(Y \mid \text{do}(X = x)) \label{eq:int}\\
  \text{Counterfactual:} &\quad P(Y_{X=x} \mid X = x', Y = y'). \label{eq:counter}
\end{align}

\noindent LLMs trained on observational text (Equation~\ref{eq:obs})
cannot in general compute interventional or counterfactual quantities
(Equations~\ref{eq:int}--\ref{eq:counter}) without explicit causal
structure.

\begin{lstlisting}[caption={Simple world model in PyTorch.}]
class WorldModel(nn.Module):
    def __init__(self, d_state, d_action, d_h):
        super().__init__()
        self.net = nn.Sequential(
            nn.Linear(d_state + d_action, d_h),
            nn.SiLU(),
            nn.Linear(d_h, d_h),
            nn.SiLU(),
            nn.Linear(d_h, d_state))

    def forward(self, s, a):
        # s: (B, d_state), a: (B, d_action)
        return self.net(torch.cat([s, a], -1))

# Eq. (8): MSE loss in latent space
loss_fn = nn.MSELoss()
\end{lstlisting}

\subsection{Emergent World Models in LLMs}

\citet{li2023emergent} showed that a transformer trained on Othello move
sequences develops an internal representation of the board state that
can be recovered by linear probing.  This suggests LLMs do develop
\emph{implicit} world models — but whether these are genuinely
causal or merely correlational is an open question.

% ---------------------------------------------------------------
\section{Open Question 3: Reasoning and Computational Depth}
\label{sec:reasoning}
% ---------------------------------------------------------------

\subsection{Transformers as Bounded-Depth Circuits}

A transformer with $L$ layers computes a function of fixed depth:

\begin{equation}
  f_\theta = f_L \circ f_{L-1} \circ \cdots \circ f_1,
  \quad \text{depth}(f_\theta) = O(L).
  \label{eq:depth}
\end{equation}

\noindent \citet{merrill2023parallelism} proved that log-precision
transformers cannot solve problems outside $\text{TC}^0$ in a single
forward pass.  Problems requiring $\omega(L)$ sequential steps therefore
cannot be solved by a single forward pass.

\subsection{Chain-of-Thought as External Memory}

Generating $T$ additional output tokens extends effective depth to
$O(L \cdot T)$, explaining why chain-of-thought (CoT) prompting
\citep{wei2022chain} dramatically improves multi-step reasoning.

The per-token log probability under CoT is

\begin{equation}
  \log P_\theta(y_t \mid r_1, \ldots, r_k, y_{<t})
  \label{eq:cot}
\end{equation}

\noindent where $r_1, \ldots, r_k$ are reasoning tokens.  The scratchpad
allows the model to re-read intermediate computations as context.

\subsection{Scaling Laws}

\citet{hoffmann2022chinchilla} derived compute-optimal scaling laws:

\begin{align}
  \mathcal{L}(N, D) &\approx E + \frac{A}{N^{\alpha}} + \frac{B}{D^{\beta}}
  \label{eq:scaling}
\end{align}

\noindent where $N$ is parameter count, $D$ is data tokens, and
$\alpha, \beta \approx 0.34$.  The irreducible loss $E$ represents the
entropy of natural language.

\begin{remark}
Emergent abilities \citep{wei2022emergent} appear as sharp transitions in
performance at certain scale thresholds.  Whether these are genuine
phase transitions or artifacts of the evaluation metric is debated.
\end{remark}

% ---------------------------------------------------------------
\section{Open Question 4: JEPA and Representation Space Prediction}
\label{sec:jepa}
% ---------------------------------------------------------------

\subsection{The JEPA Objective}

The \textbf{Joint Embedding Predictive Architecture} (JEPA)
\citep{lecun2022path,assran2023ijepa} replaces next-token prediction
with prediction in latent representation space.  Given context $x$
and target $y$, JEPA minimises

\begin{align}
  s_x &= E_x(x), \quad s_y = E_y(y) \label{eq:jepa_enc}\\
  \hat{s}_y &= g(s_x, z) \label{eq:jepa_pred}\\
  \mathcal{L}_{\text{JEPA}} &= \|g(s_x, z) - \text{sg}(s_y)\|_2^2,
  \label{eq:jepa_loss}
\end{align}

\noindent where $\text{sg}(\cdot)$ is the stop-gradient operator,
$E_y$ is an exponential moving average (EMA) of $E_x$, and $z$ is
a latent variable encoding information about $y$ not present in $x$.

\subsection{EMA Update and Collapse Prevention}

The EMA update

\begin{equation}
  \bar{\theta} \leftarrow \tau \bar{\theta} + (1 - \tau) \theta,
  \quad \tau \in [0.996, 1.0)
  \label{eq:ema}
\end{equation}

\noindent provides a slowly-evolving, stable prediction target, preventing
representational collapse without requiring contrastive negative pairs.

\begin{lstlisting}[caption={JEPA with EMA target encoder.}]
import copy

class JEPA(nn.Module):
    def __init__(self, enc, pred, tau=0.996):
        super().__init__()
        self.enc_x = enc
        self.enc_y = copy.deepcopy(enc)
        for p in self.enc_y.parameters():
            p.requires_grad_(False)
        self.predictor = pred
        self.tau = tau

    @torch.no_grad()
    def ema_update(self):
        # Eq. (14)
        for po, pe in zip(
                self.enc_x.parameters(),
                self.enc_y.parameters()):
            pe.data.mul_(self.tau).add_(
                po.data, alpha=1 - self.tau)

    def forward(self, x, y, z=None):
        s_x = self.enc_x(x)             # Eq. (12)
        with torch.no_grad():
            s_y = self.enc_y(y)         # Eq. (12)
        s_hat = (self.predictor(s_x, z)
                 if z is not None
                 else self.predictor(s_x))
        return (s_hat - s_y.detach()
                ).pow(2).mean()         # Eq. (13)
\end{lstlisting}

\subsection{VICReg: Preventing Collapse via Regularisation}

VICReg \citep{bardes2022vicreg} prevents collapse without negative pairs
through three explicit regularisation terms:

\begin{align}
  \mathcal{L}_{\text{VICReg}} &= \lambda\, \underbrace{\|z_1 - z_2\|_2^2}_{\text{invariance}}
  + \mu\, \underbrace{v(z_1, z_2)}_{\text{variance}}
  + \nu\, \underbrace{c(z_1, z_2)}_{\text{covariance}}
  \label{eq:vicreg}
\end{align}

\noindent where $v(z) = \sum_j \max(0, \gamma - \text{Std}(z_j))$ and
$c(z) = \sum_{i \neq j} [C(z)]_{ij}^2 / d$ with
$C(z) = z^\top z / (N-1)$ the normalised covariance matrix.

\subsection{H-JEPA: Hierarchical World Model}

The \textbf{Hierarchical JEPA} stacks $K$ JEPA levels at increasing
temporal scale and abstraction:

\begin{align}
  s^{(k)}_t &= \text{Pool}(s^{(k-1)}_t), \quad k = 2, \ldots, K
  \label{eq:hjepa_pool}\\
  \hat{s}^{(k)}_{t+1} &= g^{(k)}(s^{(k)}_t, a^{(k)}_t, z^{(k)})
  \label{eq:hjepa_pred}\\
  \mathcal{L}^{(k)} &= \|\hat{s}^{(k)}_{t+1} - \text{sg}(s^{(k)}_{t+1})\|_2^2.
  \label{eq:hjepa_loss}
\end{align}

\noindent Level $k=1$ makes fine-grained short-horizon predictions;
level $k=K$ makes abstract long-horizon predictions suitable for
high-level planning.

\subsection{What JEPA Does Not Solve}

\begin{remark}
JEPA substantially improves world model quality and enables gradient-based
planning (Section~\ref{sec:modular}).  However, it does \emph{not} solve
interpretability: a latent vector $[0.34, -1.82, 0.77, \ldots] \in \mathbb{R}^d$
is no more interpretable than an LLM hidden state.  The learned latent
space may be more structured, but it is still distributed and not
symbolically grounded.
\end{remark}

% ---------------------------------------------------------------
\section{Open Question 5: Modular AI and Gradient-Based Planning}
\label{sec:modular}
% ---------------------------------------------------------------

\subsection{The Six-Module Architecture}

\citet{lecun2022path} proposes replacing a monolithic LLM with six
specialised modules: Configurator, Perception, World Model (H-JEPA),
Short-Term Memory, Cost Module, and Actor.  The Actor performs
\textbf{inference-by-optimisation}:

\begin{equation}
  a^* = \arg\min_{a_{0:H}} \sum_{t=0}^{H} c(s_t, a_t),
  \quad s_{t+1} = f_\theta(s_t, a_t),
  \label{eq:planning}
\end{equation}

\noindent where the world model $f_\theta$ is fully differentiable,
so $\nabla_{a_{0:H}} \mathcal{L}$ can be computed by backpropagation —
planning without real-world trial and error.

\begin{lstlisting}[caption={Gradient-based planning through a world model.}]
def plan(world_model, cost_fn, s0,
         horizon=10, n_steps=100, lr=1e-2):
    """
    Solve Eq. (18) by gradient descent.
    world_model: (s, a) -> s_next
    cost_fn:     s      -> scalar
    """
    d_a = s0.shape[-1]  # action dim = state dim
    actions = torch.zeros(
        horizon, d_a, requires_grad=True)
    opt = torch.optim.Adam([actions], lr=lr)
    for _ in range(n_steps):
        s, total_cost = s0, 0.0
        for t in range(horizon):
            s = world_model(s, actions[t])
            total_cost = total_cost + cost_fn(s)
        opt.zero_grad()
        total_cost.backward()   # grads thru WM
        opt.step()
    return actions.detach()
\end{lstlisting}

\subsection{Information Efficiency of Self-Supervised Learning}

Self-supervised learning provides approximately $10^5$ bits of training
signal per sample (one bit per masked patch), compared to $\sim 10$ bits
per sample for supervised learning (one class label).  This
$10^4\times$ advantage makes SSL the natural training paradigm for
world models.

\subsection{Comparison of Approaches}

Table~\ref{tab:comparison} compares architectures across the five open
questions.

\begin{table}[h]
\centering
\caption{Architecture comparison across open questions.
\checkmark~= addressed, $\sim$~= partially addressed, $\times$~= unsolved.}
\label{tab:comparison}
\begin{tabular}{@{}lccccc@{}}
\toprule
\textbf{Architecture} & \textbf{Interp.} & \textbf{WM} & \textbf{Causal} & \textbf{Plan} & \textbf{Halluc.} \\
\midrule
LLM (GPT)      & $\times$ & $\sim$ & $\times$ & $\sim$ & $\times$ \\
LLM + RLHF     & $\times$ & $\sim$ & $\times$ & $\sim$ & $\sim$   \\
I-JEPA         & $\sim$   & $\sim$ & $\times$ & $\times$ & N/A    \\
H-JEPA         & $\sim$   & \checkmark & $\sim$ & \checkmark & N/A \\
H-JEPA+Modules & $\sim$   & \checkmark & $\sim$ & \checkmark & $\sim$ \\
\bottomrule
\end{tabular}
\end{table}

% ---------------------------------------------------------------
\section{An Integrated Research Agenda}
\label{sec:agenda}
% ---------------------------------------------------------------

Drawing together the five open questions, we identify a set of
high-priority research directions:

\begin{enumerate}[leftmargin=1.5em]
\item \textbf{Automated circuit discovery at scale.}
  Scaling SAEs and automated probing to $10^{12}$-parameter models
  requires new algorithmic efficiency — the bottleneck is not
  expressiveness but compute.

\item \textbf{Causal inductive biases.}
  Training objectives that explicitly reward interventional reasoning
  (Equation~\ref{eq:int}), rather than purely observational next-token
  prediction (Equation~\ref{eq:lm}).

\item \textbf{Structured latent spaces.}
  Designing latent spaces with explicit structural priors — e.g.,
  factored representations, object-centric slots — so that semantic
  dimensions are interpretable by construction.

\item \textbf{Hierarchical world models.}
  Validating H-JEPA on embodied planning tasks that require
  long-horizon reasoning across multiple abstraction levels.

\item \textbf{Evaluation benchmarks.}
  Ground-truth evaluation criteria for interpretability claims,
  causal model quality, and planning performance in open-ended
  environments.
\end{enumerate}

% ---------------------------------------------------------------
\section{Conclusion}
\label{sec:conclusion}
% ---------------------------------------------------------------

We have surveyed five grand open questions at the frontier of AI
research.  Current LLMs achieve remarkable performance but lack
interpretable circuits, explicit causal models, and robust planning
capabilities.  Architectures such as JEPA and H-JEPA address the
world model and planning questions but do not resolve interpretability.
Modular systems combining explicit world models, memory, and
gradient-based planners represent a promising direction, but require
new training algorithms, evaluation methods, and theoretical foundations.

The central message is that answering these questions requires not
only scaling existing architectures, but fundamentally reconsidering
the inductive biases of learning systems — moving from pure
next-token prediction toward architectures that explicitly represent
states, transitions, and causal structure.

% ---------------------------------------------------------------
\section*{Acknowledgements}
% ---------------------------------------------------------------

The author thanks the mechanistic interpretability and self-supervised
learning communities for making their code and models publicly available.

% ---------------------------------------------------------------
% References
% ---------------------------------------------------------------
\bibliographystyle{plainnat}
\begin{thebibliography}{99}

\bibitem[Anthropic(2024)]{anthropic2024claude}
Anthropic (2024).
\newblock The Claude 3 model family: Opus, Sonnet, Haiku.
\newblock \emph{Technical Report}, Anthropic.

\bibitem[Assran et~al.(2023)]{assran2023ijepa}
Assran, M., Duval, Q., Misra, I., Bojanowski, P., Vincent, P., Rabbat, M.,
  LeCun, Y., and Ballas, N. (2023).
\newblock Self-supervised learning from images with a joint-embedding predictive
  architecture.
\newblock In \emph{CVPR 2023}.

\bibitem[Bardes et~al.(2022)]{bardes2022vicreg}
Bardes, A., Ponce, J., and LeCun, Y. (2022).
\newblock {VICReg}: Variance-invariance-covariance regularization for
  self-supervised learning.
\newblock In \emph{ICLR 2022}.

\bibitem[Bricken et~al.(2023)]{bricken2023monosemanticity}
Bricken, T., Templeton, A., Batson, J., Chen, B., Jermyn, A., Conerly, T.,
  Turner, N., Anil, C., Denison, C., Askell, A., et~al. (2023).
\newblock Towards monosemanticity: Decomposing language models with dictionary
  learning.
\newblock \emph{Transformer Circuits Thread}.

\bibitem[Elhage et~al.(2022)]{elhage2022superposition}
Elhage, N., Hume, T., Olsson, C., Schiefer, N., Henighan, T., Kravec, S.,
  Hatfield-Dodds, Z., Lasenby, R., Drain, D., Chen, C., et~al. (2022).
\newblock Toy models of superposition.
\newblock \emph{Transformer Circuits Thread}.

\bibitem[Gemini Team(2023)]{gemini2023}
Gemini Team (2023).
\newblock Gemini: A family of highly capable multimodal models.
\newblock \emph{Technical Report}, Google DeepMind.

\bibitem[Hoffmann et~al.(2022)]{hoffmann2022chinchilla}
Hoffmann, J., Borgeaud, S., Mensch, A., Buchatskaya, E., Cai, T., Rutherford,
  E., Casas, D., Hendricks, L.~A., Welbl, J., Clark, A., et~al. (2022).
\newblock Training compute-optimal large language models.
\newblock In \emph{NeurIPS 2022}.

\bibitem[LeCun(2022)]{lecun2022path}
LeCun, Y. (2022).
\newblock A path towards autonomous machine intelligence.
\newblock \emph{OpenReview}, version 0.9.2.

\bibitem[Li et~al.(2023)]{li2023emergent}
Li, K., Hopkins, A.~K., Bau, D., Vi\'{e}gas, F., Pfister, H., and Wattenberg,
  M. (2023).
\newblock Emergent world representations: Exploring a sequence model trained on
  a synthetic task.
\newblock In \emph{ICLR 2023}.

\bibitem[Merrill and Sabharwal(2023)]{merrill2023parallelism}
Merrill, W. and Sabharwal, A. (2023).
\newblock The parallelism tradeoff: Limitations of log-precision transformers.
\newblock \emph{Transactions of the Association for Computational Linguistics},
  11:531--545.

\bibitem[Nakkiran et~al.(2021)]{nakkiran2021doubledescent}
Nakkiran, P., Kaplun, G., Bansal, Y., Yang, T., Barak, B., and Sutskever, I.
  (2021).
\newblock Deep double descent: Where bigger models and more data hurt.
\newblock In \emph{ICLR 2020}.

\bibitem[Olah et~al.(2020)]{olah2020zoom}
Olah, C., Cammarata, N., Schubert, L., Goh, G., Petrov, M., and Carter, S.
  (2020).
\newblock Zoom in: An introduction to circuits.
\newblock \emph{Distill}, \url{https://distill.pub/2020/circuits/zoom-in}.

\bibitem[OpenAI(2023)]{openai2023gpt4}
OpenAI (2023).
\newblock {GPT-4} technical report.
\newblock \emph{arXiv preprint arXiv:2303.08774}.

\bibitem[Ouyang et~al.(2022)]{ouyang2022instructgpt}
Ouyang, L., Wu, J., Jiang, X., Almeida, D., Wainwright, C.~L., Mishkin, P.,
  Zhang, C., Agarwal, S., Slama, K., Ray, A., et~al. (2022).
\newblock Training language models to follow instructions with human feedback.
\newblock In \emph{NeurIPS 2022}.

\bibitem[Pearl(2009)]{pearl2009causality}
Pearl, J. (2009).
\newblock \emph{Causality: Models, Reasoning, and Inference}.
\newblock Cambridge University Press, 2nd edition.

\bibitem[Vaswani et~al.(2017)]{vaswani2017attention}
Vaswani, A., Shazeer, N., Parmar, N., Uszkoreit, J., Jones, L., Gomez,
  A.~N., Kaiser, L., and Polosukhin, I. (2017).
\newblock Attention is all you need.
\newblock In \emph{NeurIPS 2017}.

\bibitem[Wang et~al.(2022)]{wang2022ioi}
Wang, K., Variengien, A., Conmy, A., Shlegeris, B., and Steinhardt, J. (2022).
\newblock Interpretability in the wild: a circuit for indirect object
  identification in {GPT-2} small.
\newblock In \emph{ICLR 2023}.

\bibitem[Wei et~al.(2022a)]{wei2022chain}
Wei, J., Wang, X., Schuurmans, D., Bosma, M., Chi, E., Le, Q., and Zhou, D.
  (2022a).
\newblock Chain-of-thought prompting elicits reasoning in large language models.
\newblock In \emph{NeurIPS 2022}.

\bibitem[Wei et~al.(2022b)]{wei2022emergent}
Wei, J., Tay, Y., Bommasani, R., Raffel, C., Zoph, B., Borgeaud, S., Yogatama,
  D., Bosma, M., Zhou, D., Miculivicius, D., et~al. (2022b).
\newblock Emergent abilities of large language models.
\newblock \emph{Transactions on Machine Learning Research}.

\end{thebibliography}

\end{document}
